\section{SDP Hardness Families}
\label{sec:sdp-hardness}

This section gives SDP-native lower bounds rather than embedding the LP
families of \cref{sec:lp-hardness}.  Three output levels must be kept
separate.  The first family has an intrinsic, constant gap in the optimal
value.  The second has a public Schur-complement solve but makes the requested
original-coordinate state expensive to load.  At the third level, even a
constant-normalization encoding of a canonical KKT matrix, or the first
synthesis of a tangent projector, may require linear raw-input work.  All
claims use the access and output contracts of
\cref{sec:prelim-quantum-access,sec:prelim-output-contracts}.

The geometric boundary is equally important.  The first trace-normalized
family exposes the density slice of $\Splus^3$.  In the Schur family the
$3\times3$ blocks leave only $\Splus^2$ and hence are SOCP-representable;
$4\times4$ blocks are the first members whose free Schur complement is
$\Splus^3$ and has no finite second-order-cone lift.

\subsection{An intrinsic-spectrum trace-normalized family}
\label{sec:sdp-intrinsic}

Let $n_{\rm bit}\geq1$, let
\begin{equation}
 K=16n_{\rm bit},\quad P=2K+n_{\rm bit}=33n_{\rm bit},
 \quad h=\prod_{j=1}^{n_{\rm bit}}\varsigma_j,
 \label{eq:sdp-intrinsic-parameters}
\end{equation}
and put $D_s=\diag(s,1,1)$.  There are $P$ blocks
$X_0,\ldots,X_{P-1}\in\Splus^3$.  Consecutive blocks obey
\begin{equation}
 X_i=G_iX_{i-1}G_i^T.                         \label{eq:sdp-copy}
\end{equation}
The first $K-1$ and last $K$ edges have $G_i=I$; the intervening
$n_{\rm bit}$ edges have $G_i=D_{\varsigma_j}$ in input order.  Add
$\tr X_0=1$.  If $R_0=I$ and $R_i=G_i\cdots G_1$, elimination gives exactly
\begin{equation}
 X_i=R_iZR_i^T,\qquad Z\succeq0,\qquad\tr Z=1.              \label{eq:sdp-root-lift}
\end{equation}
In the Frobenius-isometric coordinates
\begin{equation}
 \operatorname{svec}(X)=(X_{11},\sqrt2X_{12},\sqrt2X_{13},
                          X_{22},\sqrt2X_{23},X_{33}),       \label{eq:sdp-svec}
\end{equation}
congruence by $D_s$ is $\diag(1,s,s,1,1,1)$.  Each copy row
has two nonzeros, the trace row has three, and each scalar column occurs in
at most two rows.  One input sign changes only the two predecessor
coefficients in the $12$ and $13$ rows.  Ordering the copy rows along the
path proves that they are independent; their kernel is the six-dimensional
root family in \eqref{eq:sdp-root-lift}, and the trace row removes one
further dimension.

For $a<b$, write $H_{ab}=e_ae_b^T+e_be_a^T$, and set
\begin{equation}
 u=\frac1{10},\quad q=\frac KP=\frac{16}{33},\quad
 \gamma=\frac uq=\frac{33}{160}.                           \label{eq:sdp-anisotropy}
\end{equation}
Use the public costs
\begin{equation}
 C_i=\begin{cases}
  3I+uH_{23}+\gamma H_{12},&0\leq i<K,\\
  3I+uH_{23},&K\leq i<K+n_{\rm bit},\\
  3I+uH_{23}+\gamma H_{13},&K+n_{\rm bit}\leq i<P,
 \end{cases}                                                \label{eq:sdp-local-costs}
\end{equation}
and consider
\begin{equation}
 \min\left\{\frac1P\sum_{i=0}^{P-1}\inner{C_i}{X_i}:
       \eqref{eq:sdp-copy},\ \tr X_0=1,\ X_i\succeq0\right\}.
 \label{eq:sdp-intrinsic-program}
\end{equation}
Because $\sqrt{u^2+\gamma^2}<1/4$, all local costs have spectrum in
$(11/4,13/4)$.  The point $X_i=I/3$ is strictly primal feasible, and zero
equality multipliers with slacks $C_i/P$ prove strict dual feasibility.

\begin{theorem}[Intrinsic qutrit parity]
\label{thm:sdp-intrinsic}
The program \eqref{eq:sdp-intrinsic-program} has optimal values
\begin{equation}
 v_+=\frac{29}{10},\qquad v_-=\frac{14}{5}.                 \label{eq:sdp-intrinsic-values}
\end{equation}
Consequently additive-$1/25$ or target-relative-$1/100$ value estimation
requires $\Omega(P)$ raw coefficient queries.  More exactly,
\begin{equation}
 Q_{\rm add}(\epsilon)=
 \begin{cases}\lceil n_{\rm bit}/2\rceil,&\epsilon<1/20,\\0,&\epsilon\geq1/20,
 \end{cases}
 \quad
 Q_{\rm rel}(\eta)=
 \begin{cases}\lceil n_{\rm bit}/2\rceil,&\eta<1/57,\\0,&\eta\geq1/57.
 \end{cases}                                                \label{eq:sdp-value-regimes}
\end{equation}
The feasible spectrahedron has no lift over any finite product of
second-order cones.
\end{theorem}

\begin{proof}
The $H_{23}$ term is invariant under every $R_i$.  The first anchor remains
in the root frame and the last acquires the final sign $h$.  Since
$q\gamma=u$, pulling the objective through \eqref{eq:sdp-root-lift} gives
\begin{equation}
 \overline C_h=3I+uA_h,\qquad
 A_h=H_{12}+H_{23}+hH_{13}.                                \label{eq:sdp-effective-cost}
\end{equation}
For $h=+1$, $A_h$ has eigenvalues $2,-1,-1$; for $h=-1$, its signed cycle
product is negative and its eigenvalues are $1,1,-2$.  Hence
\begin{equation}
 \operatorname{spec}(\overline C_+)
 =\left\{\frac{16}5,\frac{29}{10},\frac{29}{10}\right\},
 \quad
 \operatorname{spec}(\overline C_-)
 =\left\{\frac{31}{10},\frac{31}{10},\frac{14}5\right\}.  \label{eq:sdp-intrinsic-spectra}
\end{equation}
Rayleigh--Ritz on $Z\succeq0$, $\tr Z=1$, proves
\eqref{eq:sdp-intrinsic-values}.  The common trace nine rules out an
orthogonal congruence followed by nontrivial positive rescaling.  Even after
adding a scalar multiple of $I$, the simple-largest/repeated-smallest and
repeated-largest/simple-smallest multiplicities cannot be exchanged.  The
gap is therefore intrinsic to the optimization problem.

The additive intervals are disjoint exactly when $\epsilon<1/20$ and meet
at $57/20$.  The relative intervals meet exactly at
\begin{equation}
 \eta_*=\frac{v_+-v_-}{v_++v_-}=\frac1{57},
 \qquad v_+(1-\eta_*)=v_-(1+\eta_*)=\frac{812}{285}.
\end{equation}
Below either threshold, comparison with a public separating point computes
parity with bounded error.  One coefficient query uses at most one sign query,
and the approximate-degree bound recalled in \cref{sec:prelim-lower-bound-tools}
gives the bounded-error lower bound $\lceil n_{\rm bit}/2\rceil$.  Conversely, apply
Deutsch's one-query circuit to each pair and query an unpaired last sign;
this computes $h$ in exactly that many queries.  At and above the thresholds
the displayed intersection points are public zero-query answers.

Projection onto $X_0$ identifies the feasible set with
$\mathcal D_3=\{Z\succeq0:\tr Z=1\}$.  If it had a finite SOC lift, replace
each affine lift equation $Ay=b$ by $Ay=tb$ and impose $t\geq0$.  For $t>0$
the projection is $t\mathcal D_3$.  At $t=0$, a homogeneous lift direction
could be added at arbitrary scale to a fixed feasible lift; compactness of
$\mathcal D_3$ forces its matrix projection to be zero.  The homogenized
projection is exactly $\Splus^3$, contradicting Fawzi's theorem
\cite{fawzi2018-on-representing-the-positive-semidefinite}.
\end{proof}

With $r_+=(1,1,1)^T/\sqrt3$ and $r_-=(1,-1,1)^T/\sqrt3$, strictly
complementary reduced pairs are
\begin{equation}
 Z_+^*=\tfrac12(I-r_+r_+^T),\quad Z_-^*=r_-r_-^T,
 \quad \overline S_+^*=\tfrac3{10}r_+r_+^T,
 \quad \overline S_-^*=\tfrac3{10}(I-r_-r_-^T).            \label{eq:sdp-strict-complementarity}
\end{equation}
Lifting by $X_i^*=R_iZ_h^*R_i^T$ and
$S_i^*=P^{-1}R_i\overline S_h^*R_i^T$ gives a strictly complementary
product-cone optimum.  Stationarity follows because the remaining
coefficient tuple is orthogonal to every trace-zero feasible tangent and
hence lies in the equality adjoint's range.

\subsubsection{A public, condition-one Newton step}
\label{sec:sdp-intrinsic-newton}

Use
\begin{equation}
 \Phi_\mu(X)=\frac1P\sum_i\inner{C_i}{X_i}
              -\mu\sum_i\log\det X_i,\qquad
 X_i^0=I/3,\quad\mu=1/P.                                  \label{eq:sdp-public-start}
\end{equation}
At this point, in \eqref{eq:sdp-svec},
\begin{equation}
 g_i=(C_i-3I)/P,\qquad \nabla^2\Phi_\mu(X^0)=(9/P)I.       \label{eq:sdp-public-hessian}
\end{equation}
Thus both right-hand sides of
\begin{equation}
 \widetilde A_\varsigma\Delta x=0,\qquad
 (9/P)\Delta x-\widetilde A_\varsigma^T\Delta y=-g       \label{eq:sdp-kkt}
\end{equation}
are public.  For trace-zero $Y$, define the isometry
$\mathcal W_\varsigma(Y)=P^{-1/2}(R_0YR_0^T,\ldots,R_{P-1}YR_{P-1}^T)$.
Then
\begin{equation}
 \mathcal W_\varsigma^*\nabla^2\Phi_\mu(X^0)
 \mathcal W_\varsigma=(9/P)I_5.                            \label{eq:sdp-reduced-nine}
\end{equation}
In the physical root coordinate the Hessian is $9I_5$ and
\begin{equation}
 \Delta Z_h=-\frac u9A_h.                                  \label{eq:sdp-root-direction}
\end{equation}
Its decrement and the standard full-step bound are
\begin{equation}
 \lambda=\frac1{\sqrt{150}},\qquad
 \lambda_{\rm next}\leq\frac1{(\sqrt{150}-1)^2}<\frac1{126}.
 \label{eq:sdp-decrement}
\end{equation}
Indeed $\lambda^2=\inner{uA_h}{(9I)^{-1}uA_h}=1/150$.
The post-step root eigenvalues are $(14/45,31/90,31/90)$ for $h=+1$ and
$(29/90,29/90,16/45)$ for $h=-1$, so the undamped step remains positive in
every block.  The next-decrement estimate is the universal self-concordant
bound, not a claim about a paper-specific primal--dual neighborhood.

\begin{theorem}[Newton and complete-KKT state hardness]
\label{thm:sdp-newton-state}
At \eqref{eq:sdp-public-start}, preparing the normalized root or global
primal Newton-direction state to trace distance $1/100$ requires
$\Omega(P)$ raw queries.  This remains true for every exact feasible lift of
a reduced solve $Y$ satisfying
\begin{equation}
 \frac{\norm{9Y+uA_h}_F}{\norm{uA_h}_F}\leq\frac1{10}.       \label{eq:sdp-reduced-residual}
\end{equation}
It also holds for the normalized complete primal-plus-multiplier solution of
\eqref{eq:sdp-kkt} in the displayed unit row scaling and multiplier
convention.  The latter has a public parity observable of magnitude at least
$1489/2097>7/10$.
\end{theorem}

\begin{proof}
Up to global sign,
\begin{equation}
 \ket{\delta z_h}=\frac{\ket{12}+\ket{23}+h\ket{13}}{\sqrt3}.
\end{equation}
The public contraction $T=\ket{12}\!\bra{13}+\ket{13}\!\bra{12}$ has
expectation $2h/3$.  At block $i$, both swapped coordinates acquire the same
prefix sign, so the same expectation holds for
$T_{\rm glob}=\bigoplus_iT$.  The POVM $(I\pm T)/2$ succeeds with probability
$5/6$, and trace error $1/100$ leaves constant bias.

For the residual claim, \eqref{eq:sdp-reduced-residual} is exactly
$\norm{Y-\Delta Z_h}_F\leq\eta\norm{\Delta Z_h}_F$ with
$\eta\leq1/10$.  The projective sine angle from the exact ray is at most
$\eta$.  Direct minimization gives the sharp bound
\begin{equation}
 h\langle T\rangle\geq m(\eta):=\frac23-\frac{\eta^2}{3}
 -\frac{2\sqrt2}{3}\eta\sqrt{1-\eta^2},\qquad0\leq\eta\leq\frac13.
 \label{eq:sdp-sharp-decoder}
\end{equation}
To verify it, at sine angle $s$ write the only useful tangent coefficient as
$x\in[-1,1]$.  The expectation is
$\frac23(1-s^2)-s^2+\frac{2\sqrt2}{3}s\sqrt{1-s^2}x+
\frac43s^2x^2$.  For $s\leq1/3$ its unconstrained minimizer is below $-1$,
so $x=-1$, and the result decreases with $s$.  At $\eta=1/10$ it is
$199/300-\sqrt{198}/150>17/30$.  Also
$\lambda_{\min}(I/3+Y)>3/10$.

For the complete vector, put
\begin{equation}
 \alpha=16/33,\quad c=\sqrt2\gamma/P,\quad d=\sqrt2u/9,
 \quad v_e=\begin{cases}c(1-\alpha)e,&e\leq K,\\
                          c\alpha(P-e),&K<e<P.
             \end{cases}                                  \label{eq:sdp-tent}
\end{equation}
If $\tau_i$ is the prefix sign, the three primal chains are
$-d\tau_i,-dh\tau_i,-d$ in coordinates $12,13,23$.  Switching signed path
incidence to ordinary incidence shows that the only nonzero copy multipliers
are $-\tau_ev_e$ on $12$ and $h\tau_ev_{P-e}$ on $13$.  The public
within-sector swap then has expectation $2h(D-C)/(3D+2V)$, where
\begin{align}
 D&=\frac{11n_{\rm bit}}{1350},\nonumber\\
 V&=\frac{289n_{\rm bit}}{4950}+\frac{17}{158400n_{\rm bit}},\nonumber\\
 C&=\frac{577n_{\rm bit}}{9900}+\frac1{9900n_{\rm bit}}.  \label{eq:sdp-kkt-sums}
\end{align}
Its sign is $-h$.  The leading ratio is $1489/2097$; the ratio of positive
$1/n_{\rm bit}$ corrections is $16/17>1489/2097$, so finite size only
increases the magnitude.  Reversing outcome labels decodes parity.  Each
preparer would therefore compute parity and costs
$\Omega(n_{\rm bit})=\Omega(P)$ queries.
\end{proof}

The complete-state claim is coordinate-specific: rescaling equality rows
changes multiplier amplitudes.  Also $D A_-D=-A_+$ for
$D=\diag(1,-1,1)$, so direction rays are related by an input-dependent
orthogonal gauge.  The value gap is invariant; the state theorem requests
fixed public coordinates, consistent with \cref{sec:condensation}.

\subsubsection{Ambient robustness and access objects}
\label{sec:sdp-intrinsic-access}

For arbitrary $X_i\succeq0$, define
\begin{equation}
 e_0=\tr X_0-1,\quad E_i=X_i-G_iX_{i-1}G_i^T,\quad
 e_0^2+\sum_{i=1}^{P-1}\norm{E_i}_F^2\leq\epsilon^2.       \label{eq:sdp-ambient-residual}
\end{equation}
Gauging to $Z_i=R_i^TX_iR_i$ and telescoping gives
\begin{equation}
 |P^{-1}\sum_i\tr X_i-1|<\sqrt P\epsilon,\quad
 \left|P^{-1}\sum_i\inner{C_i}{X_i}-\inner{\overline C_h}{Z_0}\right|
 <\frac{13}{4}\sqrt P\epsilon.                            \label{eq:sdp-ambient-telescope}
\end{equation}
The second bound writes the difference as suffix costs against
$Z_i-Z_{i-1}$ and uses $\norm{C_i}_{\rm op}<13/4$.  If the objective is at
most $v_h+1/100$ and
$\epsilon\leq1/(400\sqrt P)$, define
\begin{equation}
 \rho(X):=\frac{\bigoplus_iX_i}{\sum_i\tr X_i}.             \label{eq:sdp-ambient-density}
\end{equation}
The trace bounds in \eqref{eq:sdp-ambient-telescope} make this density
operator well-defined.  The public contraction
\begin{equation}
 M=\frac1{2/5}\left(\bigoplus_iC_i-\frac{57}{20}I\right)
\end{equation}
satisfies $h\tr(M\rho(X))>1/16$.  Specifically the normalized cost is
$>23/8$ for $h=+1$ and $<113/40$ for $h=-1$.  Thus any procedure which,
for every input, returns a density operator $\widetilde\rho$ satisfying
$\Dtr(\widetilde\rho,\rho(X))\leq1/100$ for some PSD tuple $X$ meeting the
displayed residual and objective conditions requires $\Omega(P)$ raw
queries.  The tuple may be selected adversarially by the procedure.  The
residual order is necessary: averaging the
$n_{\rm bit}$ neighboring negative-parity optimizers against a fixed
positive-parity input gives wrong value and residual
\begin{equation}
 \frac4{3\sqrt{n_{\rm bit}}}=\Theta(P^{-1/2}).              \label{eq:sdp-wrong-parity-witness}
\end{equation}

For access objects, use the following state-pair construction.  If a real
matrix $B$ has public support and magnitudes and maximum absolute row and
column sums at most $a$, give each supported pair $(r,c)$ its own token.
Row and column states with amplitudes
$\sqrt{|B_{rc}|/a}$ and
$\operatorname{sgn}(B_{rc})\sqrt{|B_{rc}|/a}$, completed on disjoint failure
tokens, are orthonormal and have overlap $B_{rc}/a$.  Unitary completions
therefore give an exact normalization-$a$ block encoding.

Put $y'=-\Delta y$ in \eqref{eq:sdp-kkt} and apply this to
\begin{equation}
 M_\varsigma=\begin{pmatrix}(9/P)I&\widetilde A_\varsigma^T\\
                 \widetilde A_\varsigma&0\end{pmatrix}.    \label{eq:sdp-canonical-kkt}
\end{equation}
Its maximum absolute row and column sums equal three; the trace row attains
three and a primal row has sum at most $2+9/P<3$.  One query phases the four
symmetric hidden KKT tokens.  This gives a fully specified exact
normalization-three encoding, and substitution into the decoders proves
\begin{equation}
 q_{\rm BE}\geq\left\lceil\frac{n_{\rm bit}}2\right\rceil  \label{eq:sdp-kkt-call-lower}
\end{equation}
for root, global, or complete-KKT state preparation.  The claim is relative
to this completion; input-dependent junk may contain parity and must be
charged.

Let $\Pi_\varsigma$ project onto $\ker\widetilde A_\varsigma$.  Projecting
\eqref{eq:sdp-kkt} and calculating norms gives
\begin{equation}
 \norm{\Pi_\varsigma g}_2^2=\frac3{50P},\quad
 \norm g_2^2=\frac{41}{400P},\quad
 \norm{\Pi_\varsigma\ket{\widehat g}}_2^2=\frac{24}{41}. \label{eq:sdp-projector-signal}
\end{equation}
One use of an operator-error-$1/32$ normalization-one projector encoding
then leaves signed expectation at least
$16/41-\sqrt{24/41}/16-1/1024>0.341$.  Hence first synthesis costs at least
$\lceil n_{\rm bit}/2\rceil$ raw queries.  Conversely, applying all prefix
signs to the $12,13$ chains constructs the tangent isometry; conjugating the
public trace-zero root projector gives an exact encoding with
$2n_{\rm bit}$ queries and no QRAM.  Projector synthesis is therefore
$\Theta(n_{\rm bit})$; a supplied projector already performs the aggregation.

Condition one in \eqref{eq:sdp-reduced-nine} does not describe the unscaled
saddle matrix.  Switching removes all path signs without changing singular
values.  On each individual off-diagonal coordinate, the path-difference
operator has singular values $2\sin(j\pi/(2P))$, $1\leq j<P$; in particular,
this gives $s_{\min}^+(\widetilde A_\varsigma)\leq\pi/P$.

The trace-coupled diagonal subsystem supplies the missing lower bound.  Let
the three-coordinate vector $x=(x^{(1)},x^{(2)},x^{(3)})$ be orthogonal to
its kernel, and write $m_j=P^{-1}{\bf1}^Tx^{(j)}$.
Orthogonality to every constant triple whose
coefficients sum to zero implies $m_1=m_2=m_3=:m$.  Put
$x^{(j)}=m{\bf1}+z^{(j)}$, with each $z^{(j)}$ mean zero, and define
\begin{equation}
 E^2=\sum_{j=1}^3\norm{Bz^{(j)}}_2^2,\qquad
 t=\sum_{j=1}^3x^{(j)}_0,
\end{equation}
where $B$ is ordinary path incidence.  The path Poincar\'e inequality and
telescoping give
\begin{equation}
 \sum_j\norm{z^{(j)}}_2^2\leq\frac{P^2}{4}E^2,\qquad
 |z^{(j)}_0|\leq\sqrt P\norm{Bz^{(j)}}_2.
\end{equation}
Since $3m=t-\sum_jz^{(j)}_0$, it follows that
$\norm x_2^2\leq(9/4)P^2(t^2+E^2)$.  The last parenthesis is exactly the
squared norm of the diagonal subsystem applied to $x$.  Combining diagonal
and off-diagonal coordinates, and using the maximum absolute row and column
sums three and two, yields
\begin{equation}
 \frac{2}{3P}\leq s_{\min}^+(\widetilde A_\varsigma)
 \leq\frac\pi P,\qquad
 s_{\max}(\widetilde A_\varsigma)\leq\sqrt6.               \label{eq:sdp-equality-singular-bounds}
\end{equation}
The kernel gives the saddle eigenvalue $9/P$, while every positive equality
singular value $s$ gives eigenvalues
$\{9/P\pm\sqrt{(9/P)^2+4s^2}\}/2$.  The two-sided bounds in
\eqref{eq:sdp-equality-singular-bounds} therefore prove
\begin{equation}
 \kappa_2(M_\varsigma)=\Theta(P).                          \label{eq:sdp-saddle-kappa}
\end{equation}
Its scalar off-diagonal graph is nevertheless a forest: each coordinate
forms a subdivided path and the trace multiplier joins three distinct
diagonal paths at one new vertex.  Thus its treewidth is one.

\subsection{Non-SOCP Schur-block families}
\label{sec:sdp-schur}

Fix $k\geq2$, let $K=16n_{\rm bit}$ and $L=17n_{\rm bit}+1$, and set
$a_i=\varsigma_i$ for $i\leq n_{\rm bit}$ and $a_i=1$ thereafter, with
$\tau_i=\prod_{j=1}^ia_j$.  For $X_i\in\Splus^{k+1}$, put
$F_j=(e_je_{k+1}^T+e_{k+1}e_j^T)/\sqrt2$ and
$E=e_{k+1}e_{k+1}^T$.  Impose
\begin{align}
 \inner{F_1}{X_0}&=\sqrt2,&
 \inner{F_1}{X_i}-a_i\inner{F_1}{X_{i-1}}&=0,\nonumber\\
 \inner{F_j}{X_0}&=0,&
 \inner{F_j}{X_i}-\inner{F_j}{X_{i-1}}&=0\quad(2\leq j\leq k),\nonumber\\
 \inner E{X_0}&=1,&
 \inner E{X_i}-\inner E{X_{i-1}}&=0                       \label{eq:sdp-schur-chains}
\end{align}
for $1\leq i<L$.  Rows and scalar columns have incidence at most two.  Every
feasible block is
\begin{equation}
 X_i=\begin{pmatrix}A_i&\tau_i e_1\\\tau_i e_1^T&1\end{pmatrix},
 \qquad Z_i=A_i-e_1e_1^T\succeq0.                          \label{eq:sdp-schur-map}
\end{equation}
Put $\mu_0=1/L$ and use $C_0=\mu_0\diag(I_k,2)$ on every block.

\begin{theorem}[Exact Schur path and SOCP boundary]
\label{thm:sdp-schur}
The unique, strictly complementary optimum has $Z_i=0$.  For
$t=\mu/\mu_0$, the complete central path is
\begin{equation}
 X_{\tau_i}^{(k)}(t)=
 \begin{pmatrix}tI_k+e_1e_1^T&\tau_i e_1\\\tau_i e_1^T&1\end{pmatrix},
 \quad S_i(\mu)=
 \begin{pmatrix}\mu_0I_k&-\mu_0\tau_i e_1\\
 -\mu_0\tau_i e_1^T&\mu+\mu_0\end{pmatrix}.              \label{eq:sdp-schur-path}
\end{equation}
In every public Frobenius-orthonormal tangent basis,
$H_{\rm red}(\mu)=(\mu_0/\mu)^2I=t^{-2}I$.  All primal, slack, and multiplier
coordinates are bounded for $0<\mu\leq\mu_0$.  The feasible block slice has
a finite SOCP lift exactly when $k\leq2$.
\end{theorem}

\begin{proof}
Schur complementation gives \eqref{eq:sdp-schur-map}, $\det X_i=\det Z_i$,
and $\inner{C_0}{X_i}=\mu_0(\tr Z_i+3)$.  Thus $Z_i=0$ is the unique
optimum, and $\mu_0\tr Z_i-\mu\log\det Z_i$ is uniquely minimized at
$tI_k$, proving the primal formula.  Direct inversion gives the slack.
With $S=C-\mathcal A^*y$, the nonzero chain multipliers are
\begin{equation}
 \alpha_i=\sqrt2\mu_0(L-i)\tau_i,\qquad
 \beta_i=(\mu_0-\mu)(L-i).                                \label{eq:sdp-schur-multipliers}
\end{equation}
Their divergences reproduce $C-S$.  On $0<t\leq1$, primal entries are at
most two, slack entries at most $2/L$, and
$\norm\alpha_\infty\leq\sqrt2$, $\norm\beta_\infty\leq1$.  At $t=0$ the
rank-one primal and rank-$k$ displayed slack are strictly complementary.
For dual uniqueness, put $w_i=(\tau_i e_1,1)^T$, so
$X_i^{\rm opt}=w_iw_i^T$.  Zero gap for any dual optimum gives
$S_iw_i=0$.  Stationarity in the unconstrained upper-left coordinates fixes
$(S_i)_{1:k,1:k}=\mu_0I_k$; annihilating $w_i$ then forces the last column
and row to be $(-\mu_0\tau_i e_1,\mu_0)^T$, exactly the $t=0$ limit in
\eqref{eq:sdp-schur-path}.  Finally, full row rank of $\mathcal A$ makes
$\mathcal A^*$ injective, so the equality multipliers are unique as well.

Every feasible tangent is
$\Delta X_i=\left(\begin{smallmatrix}H_i&0\\0&0\end{smallmatrix}\right)$.
Using the inverse of \eqref{eq:sdp-schur-path}, whose upper-left block is
$t^{-1}I_k$ because $\tau_i^2=1$, the reduced log-det Hessian is
\[
 \sum_i\tr(X_i^{-1}\Delta X_iX_i^{-1}\Delta Y_i)
 =\frac{\mu_0^2}{\mu^2}\sum_i\inner{H_i}{K_i},
\]
proving scalarity.  The two local centers do not commute:
$[X_+^{(k)}(t),X_-^{(k)}(t)]=2t(e_{k+1}e_1^T-e_1e_{k+1}^T)$.
Finally, \eqref{eq:sdp-schur-map} is an affine bijection from $\Splus^k$.
$\Splus^2$ is a Lorentz cone, while $\Splus^3$ has no finite SOC lift; larger
$k$ contain it as a face.  This proves the exact boundary.
\end{proof}

The noncommutativity is removable by an input-dependent congruence; in the
$Z_i$ coordinates the optimization is public.  The product-cone
representation is essential: replacing it by one unrestricted
$(k+1)L$-dimensional PSD block and killing off-block entries takes
quadratically many scalar equalities.

\subsubsection{Tight loading and the accuracy frontier}
\label{sec:sdp-schur-frontier}

For $k=3$ and $0<t\leq1$, define
\begin{equation}
 \rho_\varsigma(t)=\frac{\bigoplus_iX_{\tau_i}^{(3)}(t)}{L(3t+2)}.
\end{equation}
On the suffix $T=\{n_{\rm bit}+1,\ldots,n_{\rm bit}+K\}$, every
$\tau_i=h$.  Let $O$ equal $e_1e_4^T+e_4e_1^T$ on these blocks and zero
elsewhere.  Then
\begin{equation}
 h\tr(O\rho_\varsigma(t))=\frac{2K}{L(3t+2)}
 \geq\frac{32n_{\rm bit}}{5(17n_{\rm bit}+1)}>\frac13.    \label{eq:sdp-schur-density-bias}
\end{equation}
Thus trace-distance-$1/100$ density preparation costs $\Omega(L)$ queries.

For the complete central triple
$g_\varsigma(t)=(\operatorname{svec}X,y,\operatorname{svec}S)$,
\begin{equation}
 \norm{g_\varsigma(t)}_2^2
 =L(3t^2+2t+4)+[2+(1-t)^2]\frac{(L+1)(2L+1)}{6L}
 +\frac{5+(1+t)^2}{L}<11L.                                \label{eq:sdp-schur-triple-norm}
\end{equation}
Swapping the suffix primal coordinates $1$ and $\sqrt2\tau_i$ gives a
public contraction $W$ with
$h\langle W\rangle=2\sqrt2K/\norm{g_\varsigma(t)}_2^2>1/5$.
Hence complete-triple amplitude preparation is also $\Omega(L)$ throughout
the late tail.

Here is the promised canonical normalization-two access object.  The equality
matrix has public unit-magnitude entries and maximum absolute row and column
sums two.  Give each row and column two public slots.  A degree-two row or
column prepares the equal superposition of its two entry tokens; a degree-one
row or column sends the unused slot to its own failure label, and a degree-zero
free column goes entirely to its failure label.  Distinct rows, columns, and
failure arms use disjoint labels.  Public two-slot Hadamards, reversible token
maps, and a fixed lexicographic bijection on invalid addresses give fixed
unitary completions $L_A$ and $R_A^{(0)}$, with every hidden sign set to
$+1$ in $R_A^{(0)}$ while retaining the public predecessor minus signs.  A
phase unitary $S_\varsigma$ recognizes the predecessor token in
signed transition $i$ and phases it by $\varsigma_i$, while acting as the
identity on all other entry, failure, and invalid labels.  Extending the sign
oracle by the public dummy value $\varsigma_0=1$ makes this one coherent raw
query.  Therefore
\begin{equation}
 U_A=L_A^\dagger S_\varsigma R_A^{(0)},\qquad
 \bra{0,r}U_A\ket{0,c}=\frac{(\mathcal A_\varsigma)_{rc}}2.              \label{eq:sdp-schur-canonical-be}
\end{equation}
After padding the row and column address spaces to a common public dimension,
this is an exact, fully specified normalization-two encoding, and every ordinary,
adjoint, or controlled call uses one raw query.  Thus
$q_{\rm BE}\geq n_{\rm bit}/2$ for either output.  An alternative
input-dependent completion, junk action, prefix table, or preprocessing stage
is admissible only when all queries used to construct it are charged.

For arbitrary $X_i\succeq0$, let $d_i=\sqrt2(X_i)_{14}$ and suppose
\begin{equation}
 \norm{\mathcal A_\varsigma(X)-b}_2\leq\epsilon,\qquad
 \sum_i\inner{C_0}{X_i}\leq6,\qquad
 \epsilon\leq\frac1{4\sqrt L}.                            \label{eq:sdp-schur-robust-contract}
\end{equation}
Define the trace-normalized density output of such a point by
\begin{equation}
 \rho(X):=\frac{\bigoplus_iX_i}{\sum_i\tr X_i}.             \label{eq:sdp-schur-robust-density}
\end{equation}
It is well-defined because the root residual forces a nonzero off-diagonal
entry and a PSD matrix with such an entry has positive trace.  If
$r_0=d_0-\sqrt2$ and $r_i=d_i-a_id_{i-1}$, switching the recurrence gives
$\tau_id_i=\sqrt2+\sum_{j\leq i}\tau_jr_j$.  Summing over $T$ and applying
Cauchy--Schwarz gives $h\sum_{i\in T}d_i>K$.  PSD and the objective cap give
$\sum_i\tr X_i\leq6L$.  Since $\tr(O_iX_i)=\sqrt2d_i$ on the suffix,
\begin{equation}
 h\tr(O\rho(X))
 =\frac{\sqrt2h\sum_{i\in T}d_i}{\sum_i\tr X_i}
 >\frac{\sqrt2K}{6L}>\frac15.                              \label{eq:sdp-schur-robust-bias}
\end{equation}
Therefore preparing any $\widetilde\rho$ with
$\Dtr(\widetilde\rho,\rho(X))\leq1/100$ for some blockwise PSD point meeting
\eqref{eq:sdp-schur-robust-contract} costs $\Omega(L)$ raw queries.  The
point may depend adversarially on the input; no centrality or matrix-norm
proximity is needed.  Conversely, averaging the exact centers
obtained by flipping each sign once gives the wrong suffix parity and
\begin{equation}
 \norm{\mathcal A_\varsigma(\overline X)-b}_2
 =2\sqrt{2/n_{\rm bit}}=\Theta(L^{-1/2}).                  \label{eq:sdp-schur-wrong-mixture}
\end{equation}

For precision parameters $\epsilon,\eta>0$, let
$Q_{\rm abs}(P,\epsilon)$ be the worst-case first-preparation raw-query cost,
including setup, of producing an unconditional density output
$\widetilde\rho$ within trace distance $1/100$ of $\rho(X)$ for some
blockwise PSD $X$ with at most $P$ blocks, absolute equality residual at most
$\epsilon$, and the displayed objective cap.  Define
$Q_{\rm rel}(P,\eta)$ identically with
$\norm{\mathcal A_\varsigma(X)-b}_2/\norm b_2\leq\eta$.
For a heralded procedure, its success probability and all repetition or
amplification costs are charged as in \cref{def:heralded-output}; reusable
preprocessing is charged as in \cref{def:amortized-accounting}.

\begin{theorem}[Accuracy--query frontier]
\label{thm:sdp-accuracy-frontier}
For at most $P$ blocks and the relational output contract above,
\begin{equation}
 Q_{\rm abs}(P,\epsilon)=\Omega(\min\{P,\epsilon^{-2}\}),
 \qquad Q_{\rm rel}(P,\eta)=\Omega(\min\{P,\eta^{-2}\})  \label{eq:sdp-accuracy-frontier}
\end{equation}
for $0<\epsilon,\eta\leq1/8$.  The relative residual is the global
RHS-relative residual of \cref{def:global-residual}.  Exact-$P$ public
padding preserves the theorem, and the bounds are tight for the tuned family.
\end{theorem}

\begin{proof}
For absolute error choose
$n_{\rm bit}=\lfloor\min\{(P-1)/17,1/(2048\epsilon^2)\}\rfloor$.
Apart from a constant-size range, this is
$\Theta(\min\{P,\epsilon^{-2}\})$ and ensures the last inequality in
\eqref{eq:sdp-schur-robust-contract}.  The decoder reduces parity to every
valid output.  Since $\norm b_2=\sqrt3$, replacing $2048$ by $6144$ proves
the relative claim.

For exact $P$, add $P-L$ public $4\times4$ dummy blocks fixed by ten
one-coordinate rows to $P^{-1}I_4$, and set the common cost scale and
parameter to $1/P$.  The cap becomes $6L/P+5(P-L)/P^2$, so total trace is at
most $6L+5$ and active bias remains above $1/5$.  Dummy blocks have no
tangent coordinates, have bounded multipliers and a strictly complementary
optimum, and do not change the active non-SOCP slice.  Their RHS makes
$\norm{b_{\rm pad}}_2^2<4$, so $8192$ replaces $6144$ for exact-$P$ relative
padding.  Reading all active signs gives an exact valid output in
$O(n_{\rm bit})$ queries, while \eqref{eq:sdp-schur-wrong-mixture} matches
the residual exponent.
\end{proof}

A hidden coordinate in block $i$ is the prefix $\tau_i$ and has complexity
$\Theta(\min\{i,n_{\rm bit}\})$.  A classical scan streams all blocks with
exactly $n_{\rm bit}$ queries and constant extra workspace.  Quantumly, use
\begin{equation}
 X_\tau^{(k)}(t)=t\sum_{j=1}^ke_je_j^T+w_\tau w_\tau^T,
 \qquad w_\tau=\tau e_1+e_{k+1},                           \label{eq:sdp-schur-gram}
\end{equation}
to prepare a public local purification.  A coherent prefix-phase loop queries
and unqueries each sign, using at most $2n_{\rm bit}$ queries, logarithmic
workspace, and no QRAM.  Public coordinate predicates extend it to the
complete-triple state.  Both first-preparation complexities are therefore
$\Theta(L)$.

If $q_k=(k+1)(k+2)/2$, a path bag containing two adjacent block cliques and
their $k+1$ transition multipliers proves
\begin{equation}
 \operatorname{tw}(K_{\rm KKT})\leq2q_k+k.                 \label{eq:sdp-schur-treewidth}
\end{equation}
Thus the free central solve uses zero input queries while original-coordinate
loading uses $\Theta(L)$.  This complements bounded-treewidth chordal
conversion \cite{zhang2024-complexity-of-chordal-conversion-for}.  The bound
is for first preparation: charged cached prefixes can make later centers
cheap.  It gives no lower bound for a reduced $Z$ output or the public value.

\subsection{Abstract tree-holonomy loading}
\label{sec:sdp-tree-loading}

\begin{theorem}[Tree-holonomy loading]
\label{thm:sdp-tree-loading}
Let $\mathcal T=(V,E)$ be a public rooted tree with edge signs
$\varsigma_e$ and transported labels
$\tau_v=\prod_{e\in\operatorname{path}(r,v)}\varsigma_e$.  For public
constant-dimensional $B_+,B_-\succeq0$ of common trace $T_0$, put
\begin{equation}
 \rho_\varsigma=\frac1{|V|T_0}\bigoplus_{v\in V}B_{\tau_v},
 \qquad d=\frac12\norm{B_+/T_0-B_-/T_0}_1.                 \label{eq:sdp-tree-state}
\end{equation}
If public constant-size circuits prepare local purifications, a purification
of \eqref{eq:sdp-tree-state} needs at most $4|E|$ sign queries and logarithmic
workspace, without QRAM.  If an unknown path of length $n_{\rm bit}$ is
followed by $K=\Theta(n_{\rm bit})$ public descendants, $|V|=\Theta(n_{\rm bit})$,
and $d>0$, preparation to trace distance below a constant depending only on
$d$ and $K/|V|$ needs at least $n_{\rm bit}/2$ queries on the first
preparation.
\end{theorem}

\begin{proof}
In a public depth-first order, an edge is on the root-to-$v$ path exactly when
$v$ lies in the child subtree's preorder interval.  Prepare a uniform
entangled vertex address, query every edge, toggle a label on that interval
when the sign is negative, and unquery.  Select $B_{\tau_v}/T_0$, then reverse
the loop to erase the label.  The two compute--uncompute loops use four
queries per edge and logarithmic workspace.

For the lower bound, project onto the $K$ descendants.  This has public
probability $\alpha=K/|V|$.  Conditioned on it, the local state is $B_h/T_0$,
whose Helstrom success is $(1+d)/2$.  Guessing randomly off the tail gives
$1/2+\alpha d/2$.  Preparation error below $\alpha d/4$ leaves constant
advantage, so parity costs $n_{\rm bit}/2$ queries.  A classical tree scan
matches the linear order and may restrict itself to the union of requested
root paths.
\end{proof}

\subsection{Matrix-valued group holonomy}
\label{sec:sdp-group-holonomy}

Let $U:G\to O(d)$, put $k=3d$, and define
$W=U_{g_{n_{\rm bit}}}\cdots U_{g_1}$.  Use the copied construction with
$P=33n_{\rm bit}$ and hidden transports
$\mathcal D_g=\diag(U_g,I_d,I_d)$.  Replace $3I$ in
\eqref{eq:sdp-local-costs} by $kI_k$ and replace scalar $H_{ab}$ by
$H_{ab}(I_d)$.  Elimination gives
\begin{equation}
 \overline C_W=kI_k+uA_W,\qquad
 A_W=\begin{pmatrix}0&I_d&W^T\\I_d&0&I_d\\W&I_d&0\end{pmatrix}.          \label{eq:sdp-group-cost}
\end{equation}
Conjugating $W$ by $Q$ conjugates $A_W$ by $\diag(Q,Q,Q)$.

\begin{theorem}[Holonomy spectrum]
\label{thm:sdp-holonomy-spectrum}
\begin{equation}
 \operatorname{spec}(A_W)=
 \bigcup_{e^{i\theta}\in\operatorname{spec}(W)}
 \left\{2\cos\frac{\theta+2\pi\ell}{3}:\ell=0,1,2\right\}.              \label{eq:sdp-holonomy-spectrum}
\end{equation}
Thus $v(W)=k+u\lambda_{\min}(A_W)$ is a conjugacy-class function.  For the
natural permutation representation of $S_3$, the identity, transposition,
and three-cycle values are
\begin{equation}
 \frac{89}{10},\qquad \frac{44}{5},\qquad
 9+\frac15\cos\frac{8\pi}{9}.                             \label{eq:sdp-s3-values}
\end{equation}
\end{theorem}

\begin{proof}
On a $W$-eigenvector with eigenvalue $e^{i\theta}$, the complexified operator
is $\left(\begin{smallmatrix}0&1&e^{-i\theta}\\1&0&1\\e^{i\theta}&1&0\end{smallmatrix}\right)$,
whose characteristic polynomial is $\lambda^3-3\lambda-2\cos\theta$.
Substituting $\lambda=2\cos\phi$ proves the spectrum.  The natural $S_3$
eigenvalue lists are $(1,1,1)$, $(1,1,-1)$, and
$(1,e^{2\pi i/3},e^{-2\pi i/3})$, with least connection eigenvalues
$-1,-2,2\cos(8\pi/9)$.  Identity and three-cycles agree in the abelianization
but have different values, so the unrestricted observable is genuinely
nonabelian.
\end{proof}

Signed-permutation representations keep each svec copy row at two
unit-magnitude nonzeros, though the hidden symbol may move the predecessor
position.  The natural $S_3$ example uses $9\times9$ blocks, public costs in
$(9-1/4,9+1/4)$, and is strictly primal--dual feasible.

\begin{theorem}[$S_3$ value and Newton hardness]
\label{thm:sdp-s3-hardness}
Restrict every symbol to $\{e,t\}$ for one fixed transposition.  Additive
$1/25$ or target-relative $1/200$ value estimation takes $\Omega(P)$ raw
queries.  At $X_i=I_9/9$, $\mu=1/P$, the ambient and reduced primal Hessians
have condition number one, the KKT right-hand side is public, and
\begin{equation}
 \Delta Z_W=-\frac{u}{81}A_W,\qquad
 \lambda=\frac{\sqrt2}{30},\qquad
 \lambda_{\min}(I_9/9+\Delta Z_W)\geq\frac{44}{405}.       \label{eq:sdp-s3-newton-constants}
\end{equation}
Preparing the normalized root or global primal direction, or the normalized
complete primal-plus-multiplier KKT state in the stated row scaling, to trace
distance $1/100$ requires $\Omega(P)$ raw queries.  The complete-state
observable has magnitude greater than $33/100$.
\end{theorem}

\begin{proof}
The restricted word is identity or $t$ according to parity.  The first two
values in \eqref{eq:sdp-s3-values} differ by $1/10$ and are separated at
$177/20$; their relative intervals are disjoint below $1/177$.  A sparse
position query is simulated by a constant number of bit queries, so parity
proves value hardness.

At the public start the gradient is $(C_i-9I)/P$ and the ambient Hessian is
$(81/P)I$.  Orthogonal transport makes the normalized tangent map isometric,
so the physical root Hessian is $81I$.  Since $\tr A_W=0$, the direction is
as displayed.  Also $\norm{A_W}_F^2=18$ and
$\norm{A_W}_{\rm op}\leq2$, so
$\lambda^2=2u^2/9=1/450$ and
$1/9-2u/81=44/405$.  The next decrement is
$2/(30-\sqrt2)^2<1/400$.

For identity and a transposition, the normalized $A_W$ overlap is $7/9$.
The root-direction states have trace distance $4\sqrt2/9$, so their fixed
Helstrom measurement succeeds with probability
$\tfrac12(1+4\sqrt2/9)>0.81$.  Let $J$ swap corresponding
Frobenius-isometric svec coordinates in the $(1,2)$ and $(1,3)$ connection
blocks and act as zero on every other coordinate.  It has expectation $2/3$
for identity and $2/9$ for a transposition; centering as
$T_G=(9J-4I)/13$ gives $2h/13$ on every physical block and hence on the
global direction.

For the complete state, a public basis makes $U_t=\diag(-1,1,1)$.  Two modes
are public and one is the signed qutrit mode.  With
$d_0=\sqrt2u/81$, its primal contribution is
$D=Pd_0^2=11n_{\rm bit}/109350$ and its multiplier tent has the $V,C$ in
\eqref{eq:sdp-kkt-sums}.  Swapping $12,13$ only in the signed mode gives
\begin{equation}
 |\langle T_{\rm sign}\rangle|
 =\frac{2239504n_{\rm bit}^2+3888}
        {6759216n_{\rm bit}^2+12393}>\frac{33}{100},        \label{eq:sdp-s3-complete-bias}
\end{equation}
by cross multiplication.  These fixed decoders retain positive bias at
trace error $1/100$, and parity proves all state lower bounds.
\end{proof}

For canonical KKT access, $45$ scalar coordinates lie in each block.  A
copy-row state uses amplitude $1/3$ on two tokens and a trace-row state uses
amplitude $1/3$ on nine.  The column preparation routes predecessor tokens
through the svec involution induced by $t$, which one symbol query implements
by toggling its public orbit bit.  This exactly encodes the equality dilation
with normalization nine.  Adding $(81/P)$ times the public primal-sector
projector $P_x$ by a two-term LCU gives
\begin{equation}
 \alpha_M=9+\frac{81}{P}<12,                               \label{eq:sdp-s3-alpha}
\end{equation}
and $q_{\rm BE}\geq\lceil n_{\rm bit}/2\rceil$ for the three outputs above.
The completion is fixed; other input-dependent junk must be charged.

The tangent projector again has signal probability $24/41$.  Combining it
with the global decoder gives $48/533$.  At operator error $1/64$, the signed
expectation remains
$48/533-2\sqrt{24/41}/64-1/4096>0.065$.  First synthesis is therefore
$\Omega(n_{\rm bit})$.  Accumulating each prefix svec involution and
conjugating the public trace-zero projector gives an exact $2n_{\rm bit}$
query construction, so the bound is tight.  The scalar KKT graph is a forest:
a layer permutation makes disjoint paths, and the trace multiplier joins
distinct diagonal paths without a cycle.

If setup uses $S$ queries and later iterations use $q_t$, the canonical XOR
position convention yields
\begin{equation}
 S+\sum_tq_t\geq n_{\rm bit}/4,                            \label{eq:sdp-group-ledger}
\end{equation}
improved to $n_{\rm bit}/2$ for the clean one-query convention.  Linear
setup may store the word or prefixes and make later iterations cheap.

The unrestricted objective detects nonabelian information, but every proved
$\Omega(P)$ query theorem above restricts inputs to the abelian subgroup
$\{e,t\}\cong\mathbb Z_2$.  This is parity hardness, not a bound for a
genuinely noncommuting ordered-word promise.

The scalar value is invariant under simultaneous orthogonal gauge by
construction.  The direction decoders instead use the fixed natural
permutation basis and root frame.  A common public gauge merely conjugates
both public measurements, whereas an input-dependent output gauge must be
learned or supplied and is outside the stated output contract.

\begin{openproblem}[Intrinsically nonabelian word hardness]
\label{open:sdp-nonabelian}
Prove a local-symbol query lower bound for distinguishing conjugacy classes
of a genuinely noncommuting word promise, and transfer it through
\eqref{eq:sdp-group-cost}, without restricting the alphabet to one cyclic
subgroup.
\end{openproblem}

\subsection{Historical boundary and synthesis}
\label{sec:sdp-synthesis}

Signed switching and cycle products are classical graph invariants
\cite{kadyan2023-switching-equivalence-hermitian}; group synchronization
organizes the mechanism as gauge transport and holonomy
\cite{gao2019-the-geometry-of-synchronization-problems}, and orthogonal
synchronization is routinely relaxed by SDPs
\cite{ling2022-solving-orthogonal-group-synchronization-via}.  General group
oracle classification and symmetric-oracle problems use different promises
\cite{zhandry2015-quantum-oracle-classification,
copeland2021-quantum-query-complexity-of-symmetric}.  Fawzi's theorem gives
the exact lift boundary, not a coefficient-query lower bound.

The early $2\times2$ sparse-block examples reduce to paired LPs in their
exact-scalar form; their shared-trace variant is genuinely PSD.  The first
noncommuting $3\times3$ Schur example is still SOCP-representable, which is
why $4\times4$ blocks are the honest headline.  The two-anchor precursor used
asymmetric anchors: its two effective costs, and hence its central matrices,
have different spectra and are not related by any orthogonal similarity.

Separately, the free-$\Splus^3$ parity-central-Hessian precursor gives one
small operator result.  For
\begin{equation}
 \overline C_h=I+a(H_{12}+hH_{13}),\quad a=8/33,
 \qquad\mathcal L_h(Y)=\overline C_hY\overline C_h,         \label{eq:sdp-hessian-precursor}
\end{equation}
the reduced Hessian has condition number below five, but synthesizing a
public-constant-normalization block encoding to unscaled error $1/8$ costs
at least $n_{\rm bit}/2$ raw queries.  Indeed
$\operatorname{svec}(\mathcal L_h(E_{33}))=(a^2,0,\sqrt2ha,0,0,1)^T$;
swapping its $13,33$ coordinates has signed expectation
$2\sqrt2a/(1+a^2)^2>0.61$, and error $1/8$ leaves constant bias.  This is a
first-synthesis statement: a supplied Hessian oracle reveals parity in one
call.  In this free-$\Splus^3$ family---unlike the two-anchor family---the
two Hessians are related by an input-dependent orthogonal gauge.

Linear sparse optimal-value lower bounds already follow from LP constructions,
including Theorem~8.4 of Apers--Gribling
\cite{apers2026-quantum-speedups-linear-programming-interior}, and general SDP
value bounds predate these families \cite{apeldoorn2020-quantum-sdp-solvers}.
Mori et al.'s sparse QLS bound also uses parity propagation
\cite{mori2026-sparsity-dependent-complexity-lower-bound}.  The contribution
here is the optimization-native conjunction, not parity, propagation, or the
linear exponent.

The qutrit component imported by \cref{sec:winner-take-all} is exactly
\eqref{eq:sdp-intrinsic-parameters} with the costs
\eqref{eq:sdp-local-costs}; under that import \(N=n_{\rm bit}\).  Its spectra are
\eqref{eq:sdp-intrinsic-spectra} and its value gap is $1/10$.  Nothing here
makes that section's public-start Newton state winner-revealing; trace
allocation is governed by \cref{sec:condensation}.  As in the LP families of
\cref{sec:lp-hardness}, reduced conditioning, KKT sparsity, data access, and
requested coordinates are separate resources.  On these parity promises the
only quantum improvement is Deutsch's factor of two; classical reading takes
$n_{\rm bit}$ signs, so there is no asymptotic quantum speedup.
